\documentclass[12pt]{article}

\usepackage[a4paper, total={6in, 8in}]{geometry}
\usepackage{textcomp}

\begin{document}
\setlength{\parindent}{0pt}

\def \t {\textrightarrow}
\def \v {\vspace{1em}}
\def \bi {\begin{itemize}}
\def \ei {\end{itemize}}
\def \s[#1] {\section*{#1}}
\def \ss[#1] {\subsection*{#1}}
\def \sss[#1] {\subsubsection*{#1}}

\s[Montale]
Le "Occasioni" danno nuovo input \t parla di un occasione, di un evento, di una persona \t inaspettato \\
La natura ha un ciclo suo, che segue incessantemente \t noi non lo abbiamo, noi siamo variabili e unici \\
"Ossi di seppia" è l'inizio della sua carriera \\
In quegli anni incontra a Firenze Irma Brandeis, americana ebrea e dantista, che viene in italia per conoscere la lingua e studiare dante \t lei va a firenze, e condividono la stessa passione per dante \\
La data delle "Occasioni" è del 39 \\
Brandeis è la svolta per la vita poetica e esistenziale dell autore \t "Brand" + "eis" = fuoco + ghiaccio \t è un amore profondo \\
Questa donna popola le liriche del "visiting angel" \t riporta a immagini oniriche di angelo, e quindi Dante \t con una destrutturazione novecentesca, Beatrice sta a Dante come Irma Brandeis sta a Montale \\
"Ossi di seppia" è del 1925 \t otterra anche il nobel \\
Destrutturazione perchè, viene allontanato da Irma e si consola con le visioni onirche dei "visiting angels" \t il valore escatologico con Beatrice è meno presente pero \\
Le visint angels sono messaggere di pace, di luce, perche Irma Brandeis è Clizia \t figura che si trova nelle "Metamorfosi" di Ovidio, e nome evoca il girasole \t inoltre è Clizia rappresenta anche la pace e il calore \\
2020 romanzo chimato "Io non sono clizia" \t romanzo sulle lettere tra i due \\
Clizia è la donna che consola, e ci sono diverse liriche a lei dedicate \\
La moglie è molto presente, la chiama "mosca" e "topa" \t pero è presente anche Clizia \\
Non vuol dire deificazione morbosa, ma è un amore genuino \t qua è platonismo che si configura nell'idealizzazione verso Clizia \\
La loro intimita emotiva si ritrova nella loro intimita personale \t c'è forte condivisione, sono entrambi letterati \\
Tutta quella delusione diventa storia \t lui non è apolitico alla lettera \t il suo è un disimpegno che impegna a non impegnarsi \\
L'avvampato sfasciume con un altra finezza

\v

Lui si spende molto, nel discorso per la premiazione del nobel ritiene di non essere insignito di tale titolo \t ma lo è stato come figura rappresentativa del nostro secolo \\
Alla fine il profano è pero non molto lontano dal sacro \t non ha infatti una religione affermata, ma le "visiting angels" gli arrivano in sogno e permettono di andare oltre al meccanicismo della realta \\
Per Montale è impossibile superare il limite \t l'uomo sta dentro il limite, e non puo andare oltre \\
Parla pero della moglie (nei mottetti, es. "Ho sceso dandoti il braccio") \t lei è oltre, e vede oltre lui \t quindi c'è una realta oltre \\
"Bufera e altro" \t Montale + corrosivo e incomprensibile \t ma perche è incomprensibile? \t se d'annunzio deve essere attraversato per essere superato, se ci sono richiami onomatopeici e analogici in Pascoli \\
Ma questo non senso del reale, anche nell ermetismo, che cerca di riagganciare i fili tranciati tra poeta e realta \t e lo fa con fiugre retoriche incomprensibili, non presenti in pascoli e poeti precedenti \\
Mentre Montale fa propria la lezione di Elliot \t in uno dei suoi momenti disperati voleva dare fuoco alle sue opere \t Thomas Elliot lo incita a mantenere in vita le sue opere, e crea aggancio che si legge dietro all espressione "correlative object" \\
In elliot i nessi tra le cose sono cosi lontani, quasi impercettibili \t e come si leggono i legami tra cose distanti tra loro? \t con oggetti che incarnano il sentimento, l'occasione, spinta e ci mostrano solo il mero oggetto che tace l'occazione spinta \\
"Citare l'oggetto e tacere l'occasione spinta" \t l occasione spinta non si conosce il sentimento \\
In Montale si parla di correlativo oggettivo (no analogie assolutamente) \t il sentimento nell oggetto pero si trova per chi scrive, mentre è difficile per il lettore \\
I sentimenti arrivano attraverso l'oggetto al lettore, che pero li coglie con i suoi paradigmi \t legami con Sartre perchè ??? \\
L'apice di questa incomprensione tra chi scrive e chi legge sta in questo, nel correlativo oggettivo \t il muro vuol dire preclusione, difesa, ma lui attribuisce al muro il limite \\
Ungaretti ha il dolore nel suo animo \t ma il sentimento è solo suo, e noi lo possiamo solo leggere e interpetare \t tutto questo è dovuto a Elliot che ha stile arido e arso

\v

La fase di scrittura ostica si vede gia negli "Ossi di seppia", in cui ci sono immagini inconsuete \t alcuni dicono solo dopo nella "bufera" ma non è giusto \\
La sua visione della storia: la storia non porta un messaggio, non parla di un "magistra vitae" \t non da un messaggio e non è affidata alla fatica dei singoli, e lasciano vincere i prepotenti \\
La sua visione è una testimonianza che porta a vedere la perdita di manzoni \\
Nel 1971 \t pubblica il "Diario del 71" \t il 68 ha influenza, e lo sdoganare una liberta interiore come nel 68 viene ripresa \t in questo periodo è a Milano e lavora per il "Corriere della sera" = intellettuale moderno

\ss[Non chiederci la parola]
La lirica è il manifesto del pensiero montaliano \t fa parte degli "Ossi di seppia" \t non si puo prescindere da questa lirica per comprendere montale, e da il titlo all'intera raccolta \\
La raccolta è divisa in porizioni, e la prima porzione ha lo stesso titolo della lirica \\
È l'emblema della non valenza della parola, e la carica ontologica presente nel limite leopardiano, qua è invece assenza \\
La lirica inizia gia con una negazione \t e introduce tono colloquiale che segue l'opera \\
Assenza di assertivita nel non chiederci \t che si azzera nel termine squadri, che dia contorni nitidi al nostro animo informe \t non possiamo definire cio che abbiamo dentro \\
Le lettere di fuoco \t marchio indelebile, e evoca l'irruenza del clima dell epoca \t una assertivita ostentata, e la sicurezza che si vede nel fascismo e nella figura di mussolini \\
Non si cura neanche della sua ombra, l'immagine della canicola e della liguria sono presenti \\
Prato polveroso \t verde offuscato dalla polvere che si coglie nei climi afosi \\
Il croco è un fiore di campo, fatto come un cacile stretto con delle punte \t il suo colore spicca nel prato \\
Vuole purificare il clima politico culturale che accompagna la nostra penisola \\
La parte che chiude l opera propone assertivo / imperativo con la formula, evoca il linguaggio scientifico \t la scienza non puo \\
Storta sillaba secca ramo \t chiasmo e sinestesia \\
Codesto è dispregiativo \t trionfo del nichilismo, non letto secondo la chiave nietzschiana ma nel suo profondo senso del non potere \\
Lucida affermazione del crollo delle certezze comunicative \t la visione depositaria delle certezze \\
Critici autorevoli: Salvatore Guglielmino e Erand Grosser (anche per Pirandello) \t definiscono il termine della "ontologia del negativo", a cui si aggiunge la "teologia del negativo" \t cosi si esprime il non senso, religiosità nel nulla (qua si esprime il nichilismo) \\
Tra gli ermetici ci sono slanci religiosi, nel senso di guardare in alto e di percepire un senso religioso (a prescindere delle singole religioni) \\
Non e non dell ultimo verso sono in realta in corsivo \t l autore epsrime l'ontologia del negativo \\
"Intervista del 51", foto whatsapp
\end{document}
